% Section 6. Owners: 6.1 each person for own algorithms; 6.2 split (see docs/report/outline.md).
\section{Evaluation of the Approach}
\label{sec:evaluation}

\subsection{Why the Approach Works}

\paragraph{Leases, majority voting, resource discovery} \todo{Tejas: intuition, correctness argument and complexity for each.}

\paragraph{Leader election, replication} \todo{Abhishek: intuition, correctness argument (one leader per term, no lost accepted work) and complexity for each.}

\paragraph{SWIM failure detection}
\emph{Intuition.} Instead of every member watching every other member, each member watches one other member per period, and asks a few others for a second opinion before it raises an alarm. News of the alarm then rides on traffic that already exists.

\emph{Correctness.} A failure detector is judged by completeness (every failure is eventually detected) and accuracy (live members are not declared failed). Completeness holds because the round-robin order guarantees that every live member probes a crashed member within a bounded number of periods, after which the suspicion becomes a death after the suspect timeout. Accuracy is probabilistic, as in any detector over an asynchronous network, but two mechanisms make false positives rare: a member is suspected only after its direct probe \emph{and} all $k$ indirect probes fail, and a wrongly suspected member refutes the suspicion before it becomes a death. Refutation is safe because updates are ordered by incarnation number: \textit{alive}$(i)$ overrides any state with a smaller incarnation, \textit{suspect}$(i)$ overrides \textit{alive}$(j)$ for $j \le i$, and \textit{dead}$(i)$ overrides both for $j \le i$. Every member applies the same rules, so members that hear the same updates in any order reach the same view. Two details matter in practice. A volunteer that registers again receives a new identifier, so an acknowledgement from a different identifier at the same address is treated as a failed probe of the old identifier. And a node starts with its incarnation set to the current Unix time, so that after a restart it overrules old reports of its own death.

\emph{Complexity.} With $n$ members, each member sends one probe per period, plus $k$ indirect requests only when a probe fails, so the load per member is $O(1)$ messages per period and the group load is $O(n)$. A message carries at most nine member records, so its size is also $O(1)$. An update is forwarded $3\lceil\log_2 n\rceil$ times and reaches the whole group in $O(\log n)$ periods. Each member stores $O(n)$ state. Detection takes at most $2n-1$ periods until some member probes the failed member, plus the suspect timeout; in expectation the first probe comes after a constant number of periods.

\paragraph{Tokens, session secrets and signed tasks}
\emph{Intuition.} The coordinator only talks to volunteers it knows, each request provably comes from the volunteer whose identifier it carries, and a volunteer only runs work that provably comes from a coordinator.

\emph{Correctness.} Ed25519 is existentially unforgeable under chosen-message attacks~\cite{bernstein2012ed25519}, so without the private key an attacker cannot produce a valid signature on any message the coordinator did not sign. Because the signed message is a length-prefixed encoding of the unit, workload, lease expiry and payload, the encoding is injective: two different tasks never yield the same message, so a valid signature for one task is useless for any other. Token comparison runs in constant time over every entry, which removes the timing side channel of an early-exit comparison.

Tokens alone are not enough. Gossip spreads every volunteer identifier to all members, so after registration an identifier proves nothing, and a malicious volunteer could otherwise submit results or request work in another volunteer's name, spoiling votes and lowering that volunteer's trust score. Registration therefore also returns a fresh 256-bit secret drawn from the operating system's random source, of which the leader keeps only the SHA-256 hash. Every later request must carry the secret, and the leader compares hashes in constant time. An attacker who knows only the identifier must guess the secret, with success probability $2^{-256}$ per attempt. The secret is leader-local and needs no replication: after a failover, volunteers are unknown to the new leader and register again, receiving new secrets. Without transport encryption, an attacker who can read network traffic can still copy a secret; this, like confidentiality of task data, is out of scope.

\emph{Complexity.} Signing and verifying cost $O(|p|)$ for a payload of size $|p|$, dominated by hashing, plus a constant number of curve operations. Checking a token against $m$ entries of length up to $L$ costs $O(mL)$, which is negligible because it happens once per registration. Checking a session secret costs one hash of 32 bytes and one constant-time comparison, $O(1)$ per request.

\subsection{How the System Handles the Main Requirements}

\paragraph{Fault tolerance}
SpareCycles tolerates three kinds of fault. A crashed volunteer is detected by gossip within seconds (2.4--2.9\,s with demonstration settings, Section~\ref{sec:scenario-volunteer}) and its work is released; if gossip cannot reach it, its lease expires instead. Each unit already runs on three different volunteers, and the leader never gives the same unit to the same volunteer twice, so with majority voting a single faulty or malicious volunteer is outvoted. \todo{Tejas: confirm with measurements once majority voting (\#7) is merged.} A crashed leader is to be replaced by election, after which the new leader continues from its replicated copy of the task table. \todo{Abhishek: describe and measure leader failover once \#13 and \#14 are merged.}

\paragraph{Performance, scalability} \todo{Tejas: refer to Section~\ref{sec:performance}.}

\paragraph{Consistency, concurrency} \todo{Abhishek.}

\paragraph{Security}
Unknown volunteers cannot register: in our tests a guessed token was rejected and counted, while valid tokens were logged by name. A registered volunteer cannot act under another volunteer's identifier: in our tests all four forged requests, with no secret or with the attacker's own secret, were rejected, nothing was recorded, and the victim's genuine result still completed the job. A volunteer configured with the wrong public key found a bad signature on its first task, ran nothing and stopped; gossip then declared it dead and the leader moved its unit to an honest volunteer, so the job still finished (Section~\ref{sec:testing}). With no tokens or keys configured, the system stays open so that it can be tried immediately, and each process logs which mode is active.
